\section{Conclusion}

In this report, we have explored the orientational ordering phenomena in a coarse-grained model of polycrystalline graphene using Monte Carlo simulations. Our findings indicate the presence of distinct phases characterized by different orientational correlation behaviors consistent with the KTHNY theory of two-dimensional melting.

The validation of our simulation framework against the well-known two-dimensional XY model confirms the accuracy of our numerical methods. The observed BKT transition, marked by a change from algebraic to exponential decay of the correlation function and a corresponding peak in the susceptibility, serves as a benchmark for the next analyses. The study of the polycrystalline graphene model revealed a similar transition scenario when grain-boundary interactions are considered. These results suggest that hexatic orientational order can emerge in two-dimensional polycrystalline systems and could be described within the KTHNY theoretical framework.

Future work could confirm the nature of the transitions observed and refine the estimates of critical temperatures and exponents. Furthermore, integrating random grain extension with different sizes and positions and considering elastic distortions in crystallites could provide a more realistic representation of polycrystalline graphene.